\section{Problem, scope and decision model}
The motivating task is concrete: select a name and adjacent signature lines
on a typeset page, move the group, edit supported text, undo the change, and
retain a portable project. Rewriting source at a guessed location can move
the wrong occurrence of repeated text. A syntactically valid edit can change
line breaks, neighboring content or an unrelated graphic. A successful
compiler exit therefore does not establish that the requested edit is correct.
Refusing every request avoids many corruptions but does not make a useful editor.

The provisional question is how visual edits to compiled \LaTeX\ can become
source changes while preserving unrelated content and rejecting ambiguous or
incorrect transformations. It describes the implementation's motivation, but
an evaluation needs a narrower primary question:
\begin{quote}
For a declared subset of single-page \LaTeX\ text moves, how effectively do
constrained source patches and verification after compilation reject incorrect
transformations, and what do they cost in valid-edit refusal and latency?
\end{quote}
Preservation, useful coverage and response time are separate obligations.
The report treats the source as authoritative and the compiled PDF as an
observable view, not as a complete editable document model.

\subsection{Supported operations and explicit boundaries}
The first direct-manipulation contract covers supported text units on one
physical page, including short groups and logical paragraphs. Selecting one
soft-wrapped line may expand to its whole source paragraph; the displayed
frame must reveal that expansion. Drawing displacement retains the unit's
original space in TeX flow. It is different from moving text into another
paragraph, changing its width or reordering the logical document.

Text replacement and spacing/style controls can legitimately reflow later
content. They use weaker, separate policies. Bounded image dimensions, table
column parameters, typography and other source adapters exist, but do not
amount to arbitrary visual editing of floats, table cells, mathematics or
TikZ internals. Opaque macro effects, multiple renderings of one source object,
annotations, unsupported paint or page coordinates, and cross-page groups
can require refusal or source editing.

English, Arabic and Hebrew source/display are the present language scope.
Display success does not grant editing authority over every shaped or
bidirectional fragment. Eligibility must be demonstrated for the actual
operation, source structure and engine. Other languages and CJK input methods
are later scope; no Windows input language was added for this work.

\subsection{State and acceptance}
Let $S_r$ denote the canonical set of source bytes at revision $r$, $C_r$ the
main-document/engine configuration and recorded build context, and $P_r$ the
last accepted compiled view. A requested operation $u$ includes its intended
source object or selected unit and operation-specific parameters. The system
proposes a bounded patch $\Delta$ and privately constructs
\[
S'=\Delta(S_r), \qquad P'=\operatorname{compile}(S',C_r).
\]
An eligible build is only one acceptance condition. The proposed source must
match the native operation plan; its evidence must satisfy the operation's
policy; and the project, source revision, build and cancellation identity must
still be current. A stale or failed candidate has no authority to replace a
newer canonical state. Exact bytes alone do not restore an old authority:
Undo can return to previous text while advancing the revision.

This is a checked, partial update procedure. It does not invert arbitrary
\TeX\ execution. Multiple sources can produce the same visible text, and a
source fragment may expand into several occurrences or affect distant layout.
Tests of identity, inverse patches and selected examples support particular
invariants; they are not a theorem about every possible document.

Three boundaries must be distinguished. Native commit advances canonical
source/history. Frontend publication presents an already prepared PDF and
selection geometry. Explicit Save later changes the original project files.
A final evaluation must record these separately, including any failure after
native commit but before complete presentation. A presentation error must not
be counted as a harmless refusal if source/history already advanced.

\subsection{Purpose of this report}
The report is an implementation audit, a record of bounded development
experiments and an evaluation design. It does not assume that Texloom is the
first visual \LaTeX\ editor or a publication-ready research contribution.
The proposed contribution is an empirical account of the trade-off between
preservation checks, accepted useful edits, refused valid edits and execution
cost in a declared domain. Working behavior is not weakened to improve a
novelty claim.
